\documentclass[12pt,a4paper]{article}

%% ============================================================
%%  Packages
%% ============================================================
\usepackage[margin=2.5cm]{geometry}
\usepackage{graphicx}
\usepackage[T1]{fontenc}
\usepackage[utf8]{inputenc}
\usepackage{booktabs}
\usepackage{longtable}
\usepackage{array}
\usepackage{makecell}
\usepackage{multirow}
\usepackage{amsmath}
\usepackage{amssymb}
\usepackage{enumitem}
\usepackage{parskip}
\graphicspath{{./}{figures/}}

%% ---- AVA brand style (colours, headings, boxes, listings) ------
\usepackage{ava_style}
\renewcommand{\AVAManualTitle}{Photometry Manual v2}

%% ---- Header right-hand text (manual-specific) ------------------
\fancyhead[R]{\small\textcolor{avaInkBlue}{\textsc{Photometry Manual v2}}}

%% ---- PDF metadata -----------------------------------------------
\hypersetup{
  pdftitle   = {AGN Variability Archive --- Photometry Manual v2},
  pdfauthor  = {J.V. Hernandez Santisteban},
}

\begin{document}
%%============================================================

\avatitlepage
  {Photometry Manual}
  {2.0}
  {16 November 2026}
  {Juan V.\ Hern\'{a}ndez Santisteban}
  {jvhs1@st-andrews.ac.uk}
  {School of Physics \& Astronomy, University of St Andrews,\\
   North Haugh, St Andrews, KY16~9SS, UK}

\tableofcontents
\clearpage

%%============================================================
\section{Instruments and Cameras}
\label{sec:instruments}
%%============================================================

LCO photometry is obtained with three camera families.

\begin{table}[h!]
\centering
\small
\begin{tabular}{llllrr}
\toprule
Camera family & Prefix & Telescope & Pixel scale & FoV & Filters \\
\midrule
Sinistro & fa (fl$^\dagger$) & 1\,m & 0.389\,\ensuremath{^{\prime\prime}}/px & $26'\times26'$ & $ugrizBVRI$ \\
Spectral & fs & 2\,m Faulkes & 0.304\,\ensuremath{^{\prime\prime}}/px & $10'\times10'$ & $ugrizBVRI$ \\
MuSCAT3  & ep & 2\,m Haleakala & 0.270\,\ensuremath{^{\prime\prime}}/px & $9'\times9'$ & $griz$ (simultaneous) \\
\bottomrule
\end{tabular}
\caption{LCO camera families used by AVA.
  $^\dagger$fl is the legacy Sinistro prefix; treated identically to fa.
  Pixel scales are defined in \texttt{ava.photometry.apertures.PIXEL\_SCALES}.}
\label{tab:cameras}
\end{table}

Frames are delivered at reduction level \texttt{rlevel=91} (BANZAI pipeline),
stored as compressed multi-extension FITS (\texttt{*.e91.fits.fz}).  The
science extension is \texttt{SCI}.

%%============================================================
\section{Pipeline Overview}
\label{sec:overview-phot}
%%============================================================

Each target runs in a Prefect subflow (\texttt{ava-target/<name>}) with up
to seven stages, executed strictly in order:

\begin{enumerate}
  \item \textbf{Initialise} (\texttt{initialise}) --- create working
        directories; write per-target \texttt{target.yaml} metadata.
  \item \textbf{Reference catalogue} (\texttt{reference-catalog}) --- build
        \texttt{reference\_stars.csv}: one row per comparison star keyed by
        Gaia DR3 \texttt{source\_id}.  Run once; force-rebuild with
        \texttt{-{}-force}.
  \item \textbf{Download} (\texttt{download}) --- fetch new BANZAI frames from
        the LCO archive since the last successful run.
  \item \textbf{Source extraction} (\texttt{sextract}) --- SExtractor aperture
        photometry; write \texttt{APERDIAM} header to each catalogue.
  \item \textbf{Calibration} (\texttt{calibrate}) --- cross-match to
        \texttt{reference\_stars.csv}, per-frame weighted zero point; write
        \texttt{frame\_qc}, \texttt{measurements.parquet}, and
        \texttt{agn\_lightcurve.parquet}.
  \item \textbf{Ensemble} (\texttt{ensemble}) \textbf{[NEW]} --- PyTICS
        ensemble calibration; writes \texttt{agn\_lightcurve\_pytics.parquet}
        and \texttt{<target>\_LCO\_merged\_lightcurve.csv} (the final science product).
  \item \textbf{Plotting} (\texttt{webplot}) \textbf{[NEW]} --- draw
        light-curve, sky-diagnostic, standard-star RMS, and Swift figures;
        optionally push to web server.
  \item \textbf{Inter-calibration} (\texttt{merge}) --- CALI MCMC
        telescope-to-telescope offsets (legacy, superseded by the ensemble step).
\end{enumerate}

Data flow per frame:
\begin{lstlisting}[style=shell]
raw_data/<target>/YYYYMMDD/*.e91.fits.fz             (BANZAI L91 frames)
   |
   +-> sextractor_ava/*_cat.fits                     (SExtractor catalogues)
   |      `-- APERDIAM header, MAG_APER[n_det, n_ap]
   |
   +-> reference_stars.csv                           (built once by -o)
   |
   +-> frame_qc (photometry.sqlite)                  (per-frame ZP, QC)
   +-> measurements.parquet                          (all stars + AGN, long form)
   +-> agn_lightcurve.parquet/.csv                   (per-frame calibrated AGN)
   |
   +-> pytics/<filter>/{stars,scopes,epochs}.parquet (PyTICS solution)
   +-> agn_lightcurve_pytics.parquet/.csv            (ensemble + absolute ZP)
   +-> <target>_LCO_merged_lightcurve.csv            (shareable science product)
   |
   `-> ava_plots/lco_latest_<target>*.png/pdf        (web figures)
\end{lstlisting}

The top-level flow \texttt{ava-pipeline} submits all targets concurrently
(limited by the \texttt{ava-targets: 10} global concurrency limit) using
\texttt{ConcurrentTaskRunner}.

%%============================================================
\section{Initialisation}
\label{sec:init}
%%============================================================

Task \texttt{initialise} (\texttt{ava/pipeline/tasks/initialise.py}):

\begin{enumerate}
  \item Looks the target up in \texttt{cfg.load\_targets()}; an unknown
        \texttt{data\_id} is an error (no stray directories from typos).
  \item Creates, idempotently: \texttt{cfg.phot\_dir(target)},
        \texttt{cfg.phot\_raw(target)}, \texttt{cfg.sextractor\_dir(target)},
        \texttt{cfg.merging\_dir(target)}, and \texttt{<phot\_dir>/UVOT}.
  \item Writes \texttt{<phot\_dir>/target.yaml} once, recording the
        target-list entry and the initialisation timestamp, so each working
        directory is self-describing.
\end{enumerate}

%%============================================================
\section{Reference Catalogue}
\label{sec:refcat}
%%============================================================

\subsection{Overview}

Task \texttt{reference-catalog} (\texttt{ava/pipeline/tasks/reference.py},
science in \texttt{ava/photometry/reference.py}) builds
\texttt{reference\_stars.csv} --- the per-field catalogue of comparison stars
used by the calibration step.  Every row is keyed by \texttt{star\_id} =
Gaia DR3 \texttt{source\_id}: stable across catalogue rebuilds, and the
identifier used throughout the downstream pipeline.  The AGN always has an
\texttt{is\_agn} row (\texttt{star\_id = $-1$} if Gaia DR3 has no source at
the nucleus).

\subsection{GSPC method (default)}

The default method (\texttt{photometry.reference.method: gspc}) uses the
\textbf{Gaia DR3 Synthetic Photometry Catalogue} (VizieR I/360), which provides
AB magnitudes in SDSS $ugriz$, PS1 $y$, and Johnson-Kron-Cousins $BVRI$
directly from Gaia XP spectra.  It agrees with SDSS~DR16 to 0.01--0.02\,mag
(tested on Mrk\,817 and PG\,0844+349, Nov 2026), where APASS was off by
0.05--0.10\,mag.

LCO band mapping (no colour terms applied yet; LCO vs.\ SDSS/JKC/PS1 filter
differences are of order 0.01--0.02\,mag):

\begin{center}
\small
\begin{tabular}{ll}
\toprule
LCO band & GSPC source \\
\midrule
$u$ & SDSS $u$ (AB: $u - 0.04$) \\
$g,r,i$ & SDSS $g,r,i$ (AB) \\
$z$ & SDSS $z$ (AB: $z + 0.02$) \\
$Y$ & PS1 $y$ (AB) \\
$B,V,R,I$ & JKC $B,V,R,I$ (Vega $\to$ AB via $+[-0.09, +0.02, +0.21, +0.45]$) \\
\bottomrule
\end{tabular}
\end{center}

Quality cuts (\texttt{quality\_ok}): RUWE $\leq 1.4$ and no Gaia neighbour
within 3\,arcsec that is brighter than the star by less than 2.5\,mag (such
neighbours contaminate XP spectra).  The AGN row is never marked
\texttt{quality\_ok}.

The catalogue falls back to the APASS\,+\,model method (Section~\ref{sec:apass_model})
if fewer than 20 GSPC stars have usable $g$, $r$, $i$ measurements.

\subsection{APASS + atmosphere-model method (fallback)}
\label{sec:apass_model}

When GSPC coverage is insufficient (\texttt{photometry.reference.fallback:
apass\_model}), the pipeline uses the legacy approach:

\begin{enumerate}
  \item \textbf{APASS DR10} (AAVSO web form with \texttt{mechanize}) or
        \textbf{APASS DR9} (VizieR II/336 via \texttt{astroquery}).  Set by
        \texttt{photometry.reference.apass\_source}: \texttt{aavso} or
        \texttt{vizier}.
  \item \textbf{Gaia DR3 cross-match} for \texttt{source\_id}, GSP-Phot
        parameters, and GSPC SDSS-$u$ (used to anchor the model $u$ band).
  \item \textbf{Per-star extinction}: E(B-V) from Schlafly \& Finkbeiner
        (2011; = $0.86\times$SFD) scaled by the fraction of an exponential
        dust layer of scale height $h = 125$\,pc in front of the star at
        distance $d_\star$ (from Gaia parallax):
        \begin{equation}
          E(B-V)_\star = E(B-V)_\mathrm{LOS}
          \bigl(1 - e^{-d_\star\,|\sin b| / h}\bigr).
        \end{equation}
        Stars without usable parallaxes get the full line-of-sight value.
        \emph{Improvement over the legacy code}: the old pipeline applied the
        full column to every star regardless of distance, over-reddening
        nearby dwarfs.
  \item \textbf{Synthetic LCO photometry}: fit a reddened BT-Settl atmosphere
        (single-track grid: $\log g = 4.5$, $[\mathrm{M/H}] = 0$) to each
        star's APASS $B,g,V,r,i$ and integrate it through the LCO filter
        curves to predict $u,B,g,V,r,i,z,Y$ in AB.  Uncertainties from
        1000 covariance draws.
  \item \textbf{$u$-band anchor}: the synthetic $u$ extrapolation is biased
        by $\approx +0.2$\,mag.  A single per-field offset ties it to Gaia
        GSPC $u$ (iterative $3\sigma$ clip, minimum 10 stars).
\end{enumerate}

\subsection{Reference catalogue schema}

One row per star, saved to \texttt{reference\_stars.csv}.

\begin{center}
\small
\begin{tabular}{lll}
\toprule
Column & Type & Description \\
\midrule
\texttt{star\_id} & int64 & Gaia DR3 \texttt{source\_id} ($-1$ = AGN with no Gaia source) \\
\texttt{ra}, \texttt{dec} & float & Sky position (deg) \\
\texttt{is\_agn} & bool & True for the AGN row (always present) \\
\texttt{quality\_ok} & bool & RUWE + neighbour cuts (GSPC); \texttt{fit\_ok} (apass\_model) \\
\texttt{lco\_<b>} & float & AB magnitude in LCO band $b \in \{u,B,g,V,r,i,z,Y\}$ \\
\texttt{lco\_<b>\_err} & float & Magnitude uncertainty \\
\texttt{gaia\_g}, \texttt{gaia\_bp\_rp} & float & Gaia photometry \\
\texttt{gaia\_ruwe} & float & RUWE (quality flag) \\
\texttt{gaia\_neighbour} & bool & Contaminating Gaia neighbour within 3\,arcsec \\
\texttt{ref\_method} & str & \texttt{gspc} or \texttt{apass\_model} \\
\bottomrule
\end{tabular}
\end{center}

Provenance is recorded in \texttt{reference\_stars\_meta.yaml} alongside the
catalogue.

%%============================================================
\section{Download}
\label{sec:download}
%%============================================================

Task \texttt{download} (\texttt{ava/pipeline/tasks/download.py}).

\subsection{Daily search strategy}

For each target and each kind (\texttt{phot}: BANZAI L91 Sinistro/Spectral/MuSCAT3;
\texttt{spec}: raw L0 FLOYDS spectrum/arc/flat):

\begin{enumerate}
  \item Look up the last successful update from \texttt{ava.db}
        \texttt{download\_state}; default to \texttt{initial\_lookback\_days}
        (config) on first run.
  \item Search from (last success $-$ \texttt{overlap\_days}) to tomorrow,
        across \texttt{daily\_proposals} and all LCO spellings
        (\texttt{lco\_names}) of the target.
  \item Mode \texttt{dry\_run} (default): report found / on disk / missing
        without fetching.  Mode \texttt{live}: fetch up to
        \texttt{max\_frames\_per\_run}; advance the state only if no failures
        and not capped.
\end{enumerate}

\subsection{Spelling check}

The task also queries the lcodb photometry DB for LCO OBJECT names near the
AGN (within \texttt{unlinked\_radius\_arcmin}) that are not linked to any AVA
target --- a spelling-check for missed data.  Warnings are logged to Prefect
but do not fail the run.

%%============================================================
\section{Source Extraction}
\label{sec:sextractor}
%%============================================================

\subsection{SExtractor configuration}

Task \texttt{sextract} (\texttt{ava/pipeline/tasks/sextract.py}).

For every \texttt{*.e91.fits.fz} in \texttt{cfg.phot\_raw(target)} without a
catalogue (newest first, up to \texttt{max\_frames\_per\_run}):

\begin{enumerate}
  \item Extract the SCI extension to a unique temporary FITS file (avoids
        collisions between parallel targets --- the legacy script used a fixed
        \texttt{temp\_sextractor\_file.fits}).
  \item Determine aperture diameters: \texttt{configured\_diameters\_px(camera)}
        reads the camera-specific override (if any) from config, else the
        \texttt{default.sex} series.
  \item Run SExtractor with \texttt{-GAIN} from the FITS header,
        \texttt{-CATALOG\_TYPE FITS\_1.0}, no check-images, quiet mode.
  \item Write the aperture series to the catalogue primary header as
        \texttt{APERDIAM} (e.g.\ \texttt{"3,5,8,11,13,15,18,21,24,27,30,33"}).
  \item Atomically rename the \texttt{*.fits.part} to \texttt{*\_cat.fits}.
\end{enumerate}

The astrometry.net \texttt{solve-field} call present in the legacy script has
been dropped: its output was never the file SExtractor ran on.  Frames with
a bad BANZAI WCS are caught later during catalogue matching and repaired there.

\subsection{Aperture series and the APERDIAM header}

\begin{warnbox}
\textbf{APERDIAM is the authoritative record of which physical apertures are
in a catalogue.}  Downstream tasks must read \texttt{APERDIAM} from the
catalogue header, not infer the series from the frame's MJD.  A frame
re-extracted with the new \texttt{default.sex} has 12 apertures for all
epochs, including pre-MJD-60909 observations.
\end{warnbox}

The aperture series is \emph{append-only}: new, larger apertures may be added
at the high-index end.  Existing indices never shift.

\subsection{The 2025 aperture expansion (MJD 60909)}
\label{sec:aper-expansion}

At \textbf{MJD 60909 (2025-10-18)} the aperture series was extended by
prepending two smaller diameters (3\,px and 5\,px), shifting every existing
index by $+2$:

\begin{table}[h!]
\centering
\small
\begin{tabular}{rr rrr}
\toprule
 & & \multicolumn{3}{c}{Physical diameter (arcsec)} \\
\cmidrule(lr){3-5}
Index & Diam (px) &
  \makecell{Sinistro 1\,m\\(0.389\,\ensuremath{^{\prime\prime}}/px)} &
  \makecell{Spectral 2\,m\\(0.304\,\ensuremath{^{\prime\prime}}/px)} &
  \makecell{MuSCAT3 2\,m\\(0.270\,\ensuremath{^{\prime\prime}}/px)} \\
\midrule
  0 &  3 &  1.17 &  0.91 &  0.81 \\
  1 &  5 &  1.95 &  1.52 &  1.35 \\
  2 &  8 &  3.11 &  2.43 &  2.16 \\
  3 & 11 &  4.28 &  3.34 &  2.97 \\
  4 & 13 &  5.06 &  3.95 &  3.51 \\
  5 & 15 &  5.83 &  4.56 &  4.05 \\
  6 & 18 &  7.00 &  5.47 &  4.86 \\
  7 & 21 &  8.17 &  6.38 &  5.67 \\
  8 & 24 &  9.34 &  7.30 &  6.48 \\
  \textbf{9} & \textbf{27} & \textbf{10.50} &  8.21 &  7.29 \\
 10 & 30 & 11.67 &  9.12 &  8.10 \\
 \textbf{11} & \textbf{33} & 12.84 & \textbf{10.03} & \textbf{8.91} \\
\bottomrule
\end{tabular}
\caption{New 12-aperture series (from MJD 60909).  Bold = science aperture
  for a 5\,arcsec-radius target.}
\label{tab:aper-new}
\end{table}

\begin{table}[h!]
\centering
\small
\begin{tabular}{rr rrr}
\toprule
 & & \multicolumn{3}{c}{Physical diameter (arcsec)} \\
\cmidrule(lr){3-5}
Index & Diam (px) &
  \makecell{Sinistro 1\,m\\(0.389\,\ensuremath{^{\prime\prime}}/px)} &
  \makecell{Spectral 2\,m\\(0.304\,\ensuremath{^{\prime\prime}}/px)} &
  \makecell{MuSCAT3 2\,m\\(0.270\,\ensuremath{^{\prime\prime}}/px)} \\
\midrule
  0 &  8 &  3.11 &  2.43 &  2.16 \\
  1 & 11 &  4.28 &  3.34 &  2.97 \\
  2 & 13 &  5.06 &  3.95 &  3.51 \\
  3 & 15 &  5.83 &  4.56 &  4.05 \\
  4 & 18 &  7.00 &  5.47 &  4.86 \\
  5 & 21 &  8.17 &  6.38 &  5.67 \\
  6 & 24 &  9.34 &  7.30 &  6.48 \\
  \textbf{7} & \textbf{27} & \textbf{10.50} &  8.21 &  7.29 \\
  8 & 30 & 11.67 &  9.12 &  8.10 \\
  \textbf{9} & \textbf{33} & 12.84 & \textbf{10.03} & \textbf{8.91} \\
\bottomrule
\end{tabular}
\caption{Old 10-aperture series (before MJD 60909).  Bold = science aperture.
  Old index~7 (Sinistro) = new index~9 --- same physical aperture (27\,px).}
\label{tab:aper-old}
\end{table}

\subsection{Science aperture index resolution}

The science aperture is defined by \textbf{5\,arcsec radius} (10\,arcsec
diameter), set in \texttt{photometry.science\_aperture\_radius\_arcsec} in
\texttt{ava\_config.yaml}.

\begin{warnbox}
\textbf{Never use a hardcoded integer aperture index.}  Always call
\texttt{science\_aperture\_index\_for\_catalog(camera, n\_apertures,
diameters\_px=cat.diameters\_px)} which reads from the catalogue's own
\texttt{APERDIAM} header.  The MJD-based \texttt{science\_aperture\_index(camera,
mjd)} is a secondary method for frames whose catalogues predate the
\texttt{APERDIAM} header.
\end{warnbox}

\begin{lstlisting}[style=python]
from ava.photometry.catalogs import read_catalog
from ava.photometry.apertures import science_aperture_index_for_catalog

cat = read_catalog("frame_cat.fits")
# Preferred: use APERDIAM from the catalogue header
idx = cat.science_index("fa14")   # calls science_aperture_index_for_catalog

# Legacy / fallback: use the MJD rule
from ava.photometry.apertures import science_aperture_index
idx = science_aperture_index("fa14", mjd=61000.0)   # -> 9
\end{lstlisting}

\subsection{Aperture series extension (November 2026)}
\label{sec:aper-extension}

In November 2026 four larger apertures were appended to the new series for
MuSCAT3, whose 10\,arcsec target falls between the 33\,px (8.91\,arcsec)
and 27\,px apertures:

\begin{center}
\small
\begin{tabular}{rr rrr}
\toprule
Index & Diam (px) & Sinistro (0.389\,\ensuremath{^{\prime\prime}}/px) & Spectral (0.304\,\ensuremath{^{\prime\prime}}/px) & MuSCAT3 (0.270\,\ensuremath{^{\prime\prime}}/px) \\
\midrule
 12 & 36 & 14.0 & 10.9 & 9.7 \\
 13 & 39 & 15.2 & 11.9 & \textbf{10.5} \\
 14 & 43 & 16.7 & 13.1 & 11.6 \\
 15 & 48 & 18.7 & 14.6 & 13.0 \\
\bottomrule
\end{tabular}
\end{center}

Bold: MuSCAT3 science aperture using the extended series (index~13 = 10.53\,arcsec).
Only cameras listed in \texttt{photometry.sextractor.extended\_aperture\_cameras}
(currently \texttt{ep} = MuSCAT3) may select from the appended apertures;
Sinistro and Spectral continue to pick from the original 12.

\subsection{Re-extraction on aperture series change}

If \texttt{photometry.sextractor.reextract\_on\_aperture\_change} is true,
catalogues whose \texttt{APERDIAM} header differs from the currently
configured series are added to the extraction queue.  This is off by default.

%%============================================================
\section{Catalogue Cross-matching}
\label{sec:catmatch}
%%============================================================

Module \texttt{ava.photometry.catalogs} reads \texttt{\_cat.fits} catalogues
and cross-matches them to \texttt{reference\_stars.csv}.

\subsection{Cross-matching algorithm}

\begin{enumerate}
  \item Nearest reference star within \texttt{max\_sep\_arcsec} (default
        1.5\,arcsec).
  \item The median RA/Dec offset of those matches is removed and the match
        repeated --- absorbs small WCS shifts frame by frame.
  \item For frames where fewer than \texttt{min\_matches} (default 10)
        reference stars are found, \texttt{solve\_wcs\_from\_patterns} is
        tried: \texttt{astroalign} matches asterism triangles of the brightest
        clean detections to the reference stars and fits a similarity
        transform.  Several source/reference size combinations and both image
        parities are tried; the solution matching the most stars within 1.5\,arcsec
        wins.  LCO frames typically have parity $-1$ (RA runs west).
  \item One detection per reference star (nearest wins); one star per
        detection.
\end{enumerate}

\subsection{The AGN detection}

The AGN is identified by its \texttt{is\_agn} flag in
\texttt{reference\_stars.csv} --- not by proximity to the pointing centre.
This correctly handles offset pointings such as NGC\,4395 (AGN 4.4\,arcmin
from the pointing centre), which the legacy code silently dropped.

\subsection{Output: \texttt{measurements()}}

The \texttt{measurements(cat, matched, aperture\_index)} function returns a
long-form table (one row per matched detection) with columns:
\texttt{star\_id}, \texttt{is\_agn}, \texttt{mag}, \texttt{magerr}
(at the science aperture), \texttt{sep\_arcsec}, \texttt{x}, \texttt{y},
\texttt{fwhm\_px}, \texttt{sex\_flags}, and \texttt{frame} (the base filename).

%%============================================================
\section{Photometric Calibration}
\label{sec:calibration}
%%============================================================

Modules \texttt{ava.photometry.calibration}, \texttt{ava.photometry.process},
and task \texttt{calibrate} (\texttt{ava/pipeline/tasks/calibrate.py}).
\textbf{Status: fully ported (November 2026).}

\subsection{Zero-point algorithm (\texttt{ava.photometry.calibration})}

For each frame the following equations are applied at the science aperture
index $k$ (chosen per catalogue from \texttt{APERDIAM}):

\begin{align}
  m_{i,\mathrm{inst}} &= \texttt{MAG\_APER}_k + 2.5\log_{10}(t_\mathrm{exp})
    &&\text{(per-second instrumental mag)} \label{eq:minst}\\
  d_i &= m_{\mathrm{ref},i} - m_{i,\mathrm{inst}}
    &&\text{(offset for each standard star $i$)} \\
  w_i &= \bigl(\sigma_{i,\mathrm{inst}}^2 + \sigma_{i,\mathrm{ref}}^2 + \sigma_\mathrm{floor}^2\bigr)^{-1} \\
  \mathrm{ZP} &= \frac{\sum w_i d_i}{\sum w_i}
    &&\text{(iterated, rejecting $|d_i - \mathrm{ZP}| > \kappa \cdot 1.4826\,\mathrm{MAD}$)} \\
  \sigma_\mathrm{ZP} &= \frac{1}{\sqrt{\sum w_i}} \cdot \max\!\bigl(1,\,\sqrt{\chi^2_\nu}\bigr) \\
  m_\mathrm{AGN} &= m_\mathrm{AGN,inst} + \mathrm{ZP}, \quad
  \sigma_\mathrm{AGN}^2 = \sigma_\mathrm{AGN,inst}^2 + \sigma_\mathrm{ZP}^2
\end{align}

The AGN and the standards are measured in the \emph{same} aperture, so the
differential light loss from seeing cancels.  The following corrections are
deliberately \emph{not} applied, because each was a bug or a no-op in the
legacy code:

\begin{itemize}[noitemsep]
  \item \textbf{No aperture correction} to the largest aperture --- applying it
        to the standards only made the AGN magnitude seeing-dependent.
  \item \textbf{No airmass term} --- in-field standards share the frame's
        airmass; it cancels in $d_i$.
  \item \textbf{No colour term on the AGN} --- the legacy code used Fairall~9's
        catalogue colour for every target.  The colour slope
        $\mathrm{d}d_i / \mathrm{d}(g-i)$ is recorded as a QC number only.
\end{itemize}

Configuration (\texttt{photometry.calibration} in \texttt{ava\_config.yaml}):

\begin{center}
\small
\begin{tabular}{lll}
\toprule
Key & Default & Meaning \\
\midrule
\texttt{min\_snr}       & 20 & S/N threshold for a standard; $\sigma_\mathrm{inst} < 1.0857 / \mathrm{min\_snr}$ \\
\texttt{max\_peak\_frac}& 0.8 & \texttt{FLUX\_MAX + BACKGROUND} $<$ this $\times$ \texttt{SATURATE} \\
\texttt{sigma\_floor}  & 0.01\,mag & Added in quadrature to every standard error \\
\texttt{clip\_kappa}   & 3.0 & Clipping threshold (multiples of MAD) \\
\texttt{min\_standards}& 5   & Minimum standards for a valid ZP \\
\texttt{n\_workers}    & 4   & Parallel worker processes per target \\
\bottomrule
\end{tabular}
\end{center}

\subsection{Per-frame processing (\texttt{ava.photometry.process})}

\texttt{process\_catalog(cat\_path, ref, meta, data\_id)} handles one catalogue:

\begin{enumerate}
  \item Read the BANZAI SCI header (\texttt{header\_meta}): FILTER, EXPTIME,
        MJD-OBS, AIRMASS, L1FWHM, SATURATE, NAXIS1/2, SITEID.
  \item Read the SExtractor catalogue; select the science aperture index from
        the \texttt{APERDIAM} header.
  \item Match detections to \texttt{reference\_stars.csv} (\texttt{match\_to\_reference});
        attempt WCS repair from star patterns if fewer than 10 reference stars match.
  \item Compute the zero point and AGN magnitude via \texttt{calibrate\_frame}.
  \item Compute frame QC numbers: WCS offset, FWHM, background, ellipticity, ZP scatter.
  \item Return one \texttt{frame\_qc} dict and one \texttt{measurements} table.
\end{enumerate}

After processing all frames, \texttt{requalify} recomputes QC flags using the
per-camera/filter median ZP (to flag clouds and focus failures) --- this cannot
be done per-frame since it requires the full ensemble.

\subsection{Data products (\texttt{ava.photometry.store})}
\label{sec:data-products}

\subsubsection{\texttt{frame\_qc} table (SQLite)}

One row per frame, stored in the photometry SQLite database
(\texttt{photometry.calibration.frame\_db}, defaulting to the lcodb photometry DB).
Key columns:

\begin{center}
\small
\begin{tabular}{lll}
\toprule
Column & Type & Description \\
\midrule
\texttt{frame} & TEXT PK & Archive basename (\texttt{<tel>-<cam>-<date>-<seq>-e91}) \\
\texttt{data\_id}, \texttt{filter}, \texttt{camera} & TEXT & Target, filter, camera ID \\
\texttt{mjd\_mid} & REAL & Mid-exposure MJD \\
\texttt{airmass} & REAL & True airmass (the legacy files stored airmass $-1$) \\
\texttt{seeing\_hdr} & REAL & Header \texttt{L1FWHM} (arcsec) \\
\texttt{fwhm\_arcsec} & REAL & Median FWHM of clean matched stars \\
\texttt{aperture\_index}, \texttt{aperture\_arcsec} & & Science aperture used \\
\texttt{zp}, \texttt{zp\_err}, \texttt{zp\_scatter} & REAL & Zero point and quality \\
\texttt{n\_zp} & INTEGER & Standards used in the ZP fit \\
\texttt{colour\_slope} & REAL & $\mathrm{d}d/\mathrm{d}(g-i)$: QC only \\
\texttt{qc\_flags} & TEXT & Comma-separated reasons; empty string = clean \\
\texttt{qc\_pass} & INTEGER & 1 if \texttt{qc\_flags} is empty \\
\bottomrule
\end{tabular}
\end{center}

\subsubsection{\texttt{measurements.parquet}}

One row per (frame, star) at the science aperture.  Contains the AGN row
(\texttt{is\_agn}) alongside all matched reference stars.  This is the input
to PyTICS (via \texttt{store.load\_pytics}).

\begin{center}
\small
\begin{tabular}{lll}
\toprule
Column & dtype & Description \\
\midrule
\texttt{frame} & string & Frame basename \\
\texttt{star\_id} & int64 & Gaia DR3 source\_id ($-1$ for AGN without Gaia) \\
\texttt{is\_agn} & bool & True for the AGN row \\
\texttt{mag\_inst} & float32 & Instrumental magnitude ($m + 2.5\log t_\mathrm{exp}$) \\
\texttt{magerr} & float32 & SExtractor aperture error \\
\texttt{x}, \texttt{y} & float32 & Pixel centroid \\
\texttt{sex\_flags} & int16 & SExtractor FLAGS \\
\texttt{sep\_arcsec} & float32 & Separation from reference position \\
\texttt{used\_in\_zp} & bool & Contributed to the frame's zero point \\
\bottomrule
\end{tabular}
\end{center}

\subsubsection{\texttt{agn\_lightcurve.parquet} and \texttt{agn\_lightcurve.csv}}

The calibrated AGN light curve, built from \texttt{measurements.parquet} joined
with \texttt{frame\_qc}:

\begin{equation}
  m_\mathrm{AGN} = m_\mathrm{inst,AGN} + \mathrm{ZP}, \qquad
  \sigma_\mathrm{AGN} = \sqrt{\sigma_\mathrm{inst}^2 + \sigma_\mathrm{ZP}^2}
\end{equation}

Columns include: \texttt{frame}, \texttt{mjd\_mid}, \texttt{filter},
\texttt{camera}, \texttt{tel}, \texttt{aperture\_arcsec}, \texttt{mag},
\texttt{mag\_err}, \texttt{mag\_inst}, \texttt{magerr}, \texttt{zp},
\texttt{zp\_err}, \texttt{qc\_pass}, \texttt{qc\_flags}.

\begin{notebox}
\textbf{Legacy \texttt{lco\_latest.pkl} is superseded.}  The new pipeline writes
\texttt{agn\_lightcurve.parquet} and \texttt{agn\_lightcurve.csv} instead.
\texttt{measurements.parquet} replaces the array-valued pkl columns with a
clean long-format Parquet file.
\end{notebox}

\subsection{PyTICS interface (\texttt{store.load\_pytics})}

\texttt{load\_pytics(phot\_dir, db, data\_id, filt, qc\_only)} joins
\texttt{measurements.parquet} with \texttt{frame\_qc} to produce the long table
expected by PyTICS: one row per (frame, star) with columns MJD, Filter,
star\_id, is\_agn, mag\_inst, magerr, telid, camera, airmass, seeing.

%%============================================================
\section{Light-Curve Plots}
\label{sec:plots}
%%============================================================

Task \texttt{webplot} (\texttt{ava/pipeline/tasks/webplot.py}).

\begin{warnbox}
\textbf{Status: not yet ported.}  The webplot task stub raises
\texttt{NotImplementedError}.
\end{warnbox}

\subsection{Planned behaviour}

\begin{enumerate}
  \item Load \texttt{lco\_latest.pkl}; apply quality cuts (\texttt{flag==0},
        \texttt{dist < 3.5}\,arcsec, \texttt{n\_std >= 3}).
  \item Convert calibrated magnitudes to flux density in mJy:
        \begin{equation}
          f_\nu\,[\text{mJy}] = 10^{(m_\mathrm{cal} - 8.9)/(-2.5)} \times 10^3.
        \end{equation}
  \item Write \texttt{lco\_latest\_<target>.pdf/png} (full timeline),
        \texttt{*\_last6.*} (last 60 days), and \texttt{*\_swift.*} (Swift UV/X-ray).
  \item If \texttt{paths.web\_server\_host} is set, copy PNGs over scp
        (key-based auth only).
\end{enumerate}

%%============================================================
\section{Ensemble Calibration (PyTICS)}
\label{sec:ensemble}
%%============================================================

Module \texttt{ava.photometry.ensemble} and task \texttt{ensemble}
(\texttt{ava/pipeline/tasks/ensemble.py}).
\textbf{Status: fully ported (November 2026).}

\subsection{Overview}

The ensemble step replaces the CALI/pycali MCMC inter-calibration with a
direct port of PyTICS \citep[Vieliute et~al.\ 2025, RASTI,
\texttt{doi:10.1093/rasti/rzaf021}]{Vieliute2025}, vectorised with Numba.
It operates on the \texttt{measurements.parquet} output of the calibration
step and writes the final science light curve.

\subsection{Model}

For every comparison star $s$ at epoch $t$ on scope $k$ (scope $=$ telescope
identifier $+$ camera identifier):
\begin{align}
  m_\mathrm{obs}^{s,t,k} &= m^*_s + \Delta T_t + \Delta S_k + C_k\,(c_s - c_0)
  \label{eq:pytics-model} \\
  \sigma_\mathrm{tot}^2 &= \sigma_\mathrm{phot}^2 + \sigma^{*2}_s + \sigma^2_T(t) + \sigma^2_{S,k}
\end{align}
where $m^*_s$ is the star's true mean magnitude, $\Delta T_t$ the per-epoch
transparency correction (zero-mean within each scope), $\Delta S_k$ the
scope offset (zero-mean over all data), and $C_k\,(c_s - c_0)$ the
per-scope colour term (\textbf{AVA extension}; $c = g - i$ from the reference
catalogue).

The iterative solution (up to \texttt{max\_loops = 200}) follows the PyTICS
\texttt{Filter.Corr} scheme exactly, using the \texttt{arx} extra-variance
optimal average (K.\ Horne, implemented as a Numba JIT kernel).  Convergence
is declared when every $|\Delta p|/\sigma_p < \texttt{tiny} = 10^{-4}$ or
the gauge-free solution moves by $<10^{-5}$\,mag for three consecutive
iterations.

The AGN is \textbf{always excluded} from the comparison stars.

\subsection{Per-scope colour terms}

After a full fit the ensemble optionally re-fits with per-scope colour terms
$C_k$ and runs \texttt{colour\_term\_check}: two reference-camera metrics
(camera-to-camera additive offset RMS and same-night flux-ratio scatter) are
compared with and without the terms.  The terms are kept for a filter only
if the mean improvement $\geq \texttt{colour\_check\_min\_gain} = 5\%$.
The decision is stored in \texttt{meta['colour\_check']} and reused by daily
incremental updates.

\subsection{Incremental update}

For the daily run (\texttt{mode = "auto"}), new frames are added to the stored
solution in tens of milliseconds per filter (\texttt{ensemble.update}): one
\texttt{arx} call per new epoch, keeping star means and scope offsets fixed.
A full re-fit is triggered automatically when: a new camera appears, more than
10\% of all epochs come from updates, or new-epoch $\sigma_T$ exceeds
$2\times$ the full-fit median.

\subsection{Absolute zero point}

After the full fit, \texttt{absolute\_zp} maps the ensemble star means onto
\texttt{reference\_stars.csv}:
\begin{equation}
  m_\mathrm{ref} - m^*_s = \mathrm{ZP} + C\,(c_s - c_0), \qquad c = g - i
\end{equation}
weighted, iteratively $3\sigma$-clipped over $\sim 100$ stars.  The AGN
colour $c_\mathrm{AGN}$ is determined from its own calibrated magnitudes
(iterated), so the zero point is evaluated at the AGN colour.

\subsection{Flux-flux intercalibration}

After ensemble calibration, \texttt{ava.photometry.intercal} fits a linear
flux model per scope against the Sinistro reference:
\begin{equation}
  F_\mathrm{scope}(t) = a \cdot F_\mathrm{ref}(t) + b
\end{equation}
where $b$ captures the \emph{additive} host-galaxy flux difference from
aperture mismatch (MuSCAT3 at 8.91\,arcsec captures less host than Sinistro
at 10.5\,arcsec).  The stars fixed $a \approx 1$ already; $a = 1$ is enforced
by the additive-only mode.  The corrected flux is $F' = F - b$.
Configured cameras for correction: \texttt{intercal\_correct\_cameras: [ep]}
(MuSCAT3 only).

\subsection{Data products}
\label{sec:ensemble-products}

\begin{description}
  \item[\texttt{pytics/<filter>/{stars,scopes,epochs,star\_scope}.parquet}]
    The stored solution.  Reloaded for incremental updates and for
    \texttt{ensemble.apply}.
  \item[\texttt{pytics/absolute\_zp\_stars.parquet}]
    Per-star $m_\mathrm{ref} - m^*$ for all filters (QC on the absolute ZP fit).
  \item[\texttt{pytics/intercal.csv}]
    Flux-flux fit table: scope, $a$, $b$, $b_1$, $\chi^2_\nu$, \texttt{corrected}.
  \item[\texttt{agn\_lightcurve\_pytics.parquet/.csv}]
    Ensemble + absolute ZP + intercalibrated AGN light curve.
    Key columns: \texttt{MJD}, \texttt{Filter}, \texttt{scope}, \texttt{mag},
    \texttt{mag\_err}, \texttt{mag\_err\_rand}, \texttt{mag\_ic}, \texttt{mag\_ic\_err},
    \texttt{F\_ic} (mJy), \texttt{DMAGT}, \texttt{DMAGS}, \texttt{CKterm},
    \texttt{rms\_t}, \texttt{rms\_sc}, \texttt{n\_stars}, \texttt{good}.
  \item[\texttt{<target>\_LCO\_merged\_lightcurve.csv}]
    The \textbf{shareable science product}: \texttt{good} epochs only, in mJy
    and AB mag with separate total and epoch-to-epoch (random) errors.
    Copied to \texttt{plots\_dir}.
\end{description}

\begin{notebox}
\textbf{Group by scope = telescope + camera.}  Sinistro detectors (fa-series)
move between telescopes.  The \texttt{scope} column combines both so each
camera is treated as a single comparison-star ensemble regardless of which
mount it was on.
\end{notebox}

\begin{lstlisting}[style=python]
# Manual run for one target (full refit)
from ava.photometry.ensemble import run_target
from ava.config import cfg
report = run_target(cfg.phot_dir("fairall_9"), "data/lcodb/photometry.sqlite",
                    "fairall_9", filters=["gp", "rp", "ip"], mode="full")

# Read the final merged light curve
import pandas as pd
lc = pd.read_parquet("/data3/jvhs1/agn/fairall_9/agn_lightcurve_pytics.parquet")
good = lc[lc["good"]]
print(good[["MJD","Filter","mag_ic","mag_ic_err","F_ic","scope"]].head())
\end{lstlisting}

%%============================================================
\section{Plotting}
\label{sec:plotting}
%%============================================================

Task \texttt{webplot} (\texttt{ava/pipeline/tasks/webplot.py}).
\textbf{Status: fully ported (November 2026).}

The webplot task calls \texttt{ava.plots.webplots} and writes to
\texttt{cfg.plots\_dir}:

\begin{description}
  \item[\texttt{lco\_latest\_<target>.png/pdf}] Full time-baseline light curve
    per filter (mJy vs MJD), FLOYDS epoch ticks on the top panel, last
    photometry and spectroscopy dates in the title.
  \item[\texttt{lco\_latest\_<target>\_last6.png/pdf}] Last
    \texttt{plots.recent\_days} (default 60) days, y-axis auto-scaled.
  \item[\texttt{lco\_latest\_<target>\_<filter>\_sky.png/pdf}] Seeing,
    airmass, and zero point vs time per filter, one series per site.
  \item[\texttt{lco\_latest\_<target>\_stan.png/pdf}] Magnitude vs median
    RMS for all field stars; AGN overlaid as a diamond.  Diagnostic for
    photometric quality: the AGN sits above the stellar sequence when
    variability exceeds the photon-noise floor.
  \item[\texttt{lco\_latest\_<target>\_swift.png/pdf}] Swift UV + X-ray
    light curves (if the target is in \texttt{paths.swift\_db}).
\end{description}

Spectroscopic epoch ticks are read directly from the lcodb spectra SQLite DB
(FLOYDS science spectra, \texttt{rlevel=0}), falling back to the last FLOYDS
date in \texttt{ava.db}.  If \texttt{paths.web\_server\_host} is set, the
PNGs are pushed to the web server via scp (key-based auth only).

Configuration (\texttt{plots} section of \texttt{ava\_config.yaml}):

\begin{center}
\small
\begin{tabular}{lll}
\toprule
Key & Default & Meaning \\
\midrule
\texttt{filter\_dict} & \texttt{data/reference/filter\_dict.txt} & filter names, wavelengths, colours \\
\texttt{recent\_days} & 60 & window for the ``last6'' plot \\
\texttt{rms\_last\_epochs} & 500 & epochs per star in the standards RMS plot \\
\texttt{formats} & \texttt{[png, pdf]} & output formats \\
\bottomrule
\end{tabular}
\end{center}

%%============================================================
\section{Inter-calibration (Legacy CALI)}
\label{sec:intercal}
%%============================================================

Task \texttt{merge} (\texttt{ava/pipeline/tasks/merge.py}).

\begin{warnbox}
\textbf{Status: superseded.}  The \texttt{merge} task stub remains for
backward compatibility but raises \texttt{NotImplementedError}.  The ensemble
step (Section~\ref{sec:ensemble}) replaces CALI/pycali for all new data.
\end{warnbox}

The legacy CALI MCMC inter-calibration (\texttt{intercalibration.mcmc\_steps = 10000})
is documented in Appendix~\ref{app:legacy} for reference.

%%============================================================
\section{Quality Control}
\label{sec:qc}
%%============================================================

\begin{table}[h!]
\centering
\small
\begin{tabular}{cl}
\toprule
\texttt{flag} & Meaning \\
\midrule
0 & Good frame (\texttt{n\_std $\geq$ 3} and seeing $<$ 9\,arcsec) \\
1 & Flagged: too few standards or seeing above limit \\
\bottomrule
\end{tabular}
\end{table}

QC flag reasons stored in \texttt{frame\_qc.qc\_flags}
(comma-separated; empty = clean):

\begin{center}
\small
\begin{tabular}{ll}
\toprule
Flag & Condition \\
\midrule
\texttt{no\_zp}       & ZP fit failed (too few standards or non-finite result) \\
\texttt{few\_matched} & fewer than \texttt{min\_matched} (30) reference stars matched \\
\texttt{agn\_missing} & AGN not detected in this frame \\
\texttt{agn\_flagged} & AGN detected with non-zero SExtractor FLAGS \\
\texttt{zp\_scatter}  & star-to-star ZP scatter $>$ \texttt{max\_zp\_scatter} (0.05\,mag) \\
\texttt{seeing}       & median FWHM $>$ \texttt{max\_fwhm\_arcsec} (4.5\,arcsec) \\
\texttt{zp\_offset}   & $|\mathrm{ZP} - \mathrm{med}(\mathrm{ZP}_\mathrm{camera/filter})|$ $>$ 0.5\,mag (clouds) \\
\bottomrule
\end{tabular}
\end{center}

Per-filter overrides are applied for $u'$ (fewer standards expected) and $z_s, Y$
(larger scatter until LCO--SDSS-$z$/PS1-$y$ colour terms exist):

\begin{center}
\small
\begin{tabular}{lll}
\toprule
Filter & Override key & Value \\
\midrule
$u'$ & \texttt{min\_matched} & 10 \\
$u'$ & \texttt{max\_zp\_scatter} & 0.08\,mag \\
$z_s$ & \texttt{max\_zp\_scatter} & 0.12\,mag \\
$Y$ & \texttt{max\_zp\_scatter} & 0.25\,mag \\
\bottomrule
\end{tabular}
\end{center}

%%============================================================
\section{Known Limitations and AVA v2 Status}
\label{sec:limitations}
%%============================================================

\begin{enumerate}

\item \textbf{Calibration task fully ported (resolved, Nov 2026).}
      Science in \texttt{ava.photometry.calibration} and \texttt{ava.photometry.process};
      outputs in \texttt{ava.photometry.store}.

\item \textbf{Pickle superseded (resolved, Nov 2026).}  \texttt{lco\_latest.pkl}
      replaced by \texttt{agn\_lightcurve.parquet/.csv} and
      \texttt{measurements.parquet}.

\item \textbf{Ensemble and webplot tasks fully ported (resolved, Nov 2026).}
      PyTICS ensemble in \texttt{ava.photometry.ensemble}; intercalibration in
      \texttt{ava.photometry.intercal}; plotting in \texttt{ava.plots.webplots}.

\item \textbf{Merge task superseded (resolved).}  The old CALI/pycali merge
      task is replaced by the ensemble step.  The \texttt{merge} stub remains
      but is not called by the daily pipeline.

\item \textbf{SExtractor as external binary (open).}  Replace with the
      \texttt{sep} Python package or \texttt{photutils}.  The sextract task
      already isolates the binary call in \texttt{extract\_one()} to ease the
      swap.

\item \textbf{pycali\_jvhs undocumented (open).}  Private package; document
      install steps or replace with PyROA.

\item \textbf{agn\_field\_id superseded (resolved).}  The old pipeline
      identified the AGN by its row index in \texttt{apass\_pytics.csv}.
      AVA~v2 uses \texttt{is\_agn} in \texttt{reference\_stars.csv}, keyed
      by Gaia DR3 \texttt{source\_id} --- robust to catalogue rebuilds.

\item \textbf{Legacy MJD typo for aperture expansion (resolved).}
      The old script compared against \texttt{mjd > 66909} (instead of 60909).
      The new code reads \texttt{APERDIAM} from the catalogue header; the MJD
      rule is a secondary fallback only.

\end{enumerate}

%%============================================================
\appendix
\section{Useful Commands}
\label{app:commands}
%%============================================================

\begin{lstlisting}[style=shell]
# Initialise + build reference catalogue for one target
python scripts/run_pipeline.py -o --target fairall_9

# Download new frames (dry run, then live)
python scripts/run_pipeline.py -d --target fairall_9
python scripts/run_pipeline.py -d --target fairall_9 --mode live

# Run SExtractor on new frames
python scripts/run_pipeline.py -e --target fairall_9

# Run calibration (new frames only)
python scripts/run_pipeline.py -c --target fairall_9

# Force full recalibration
python scripts/run_pipeline.py -c --target fairall_9 --force

# Inspect the reference catalogue
python -c "
import pandas as pd
df = pd.read_csv('/data3/jvhs1/agn/fairall_9/reference_stars.csv')
print(df[df.is_agn])
print(df[df.quality_ok][['star_id','lco_g','lco_r','lco_i','gaia_ruwe']].describe())
"

# Read the calibrated AGN light curve
python -c "
import pandas as pd
lc = pd.read_parquet('/data3/jvhs1/agn/fairall_9/agn_lightcurve.parquet')
print(lc[lc.qc_pass == 1][['mjd_mid','filter','mag','mag_err','zp','camera']].head(10))
"

# Read frame QC for a target / check failure modes
python -c "
from ava.photometry.store import read_frame_qc
fq = read_frame_qc('data/lcodb/photometry.sqlite', 'fairall_9')
print(fq.groupby('qc_flags').size().sort_values(ascending=False).head(10))
"

# Load PyTICS long table (all stars, QC-passing frames, one filter)
python -c "
from ava.photometry.store import load_pytics
t = load_pytics('/data3/jvhs1/agn/fairall_9',
                'data/lcodb/photometry.sqlite', 'fairall_9', filt='gp')
print(t[t.is_agn][['MJD','mag_inst','magerr','zp','camera']].head())
"

# Recompute QC flags after changing thresholds (no recalibration)
python -c "
from ava.photometry.process import requalify
from ava.config import cfg
fq = requalify('data/lcodb/photometry.sqlite', 'fairall_9',
               dict(vars(cfg.photometry.qc)))
print(fq['qc_pass'].value_counts())
"

# Audit aperture index
python -c "
from ava.photometry.apertures import audit_aperture_index
print(audit_aperture_index(7, 'fa14', 61000.0))
# -> {'stored_index': 7, 'stored_diameter_arcsec': 8.17,
#     'correct_index': 9, 'correct_diameter_arcsec': 10.50, 'ok': False}
"

# Print the full aperture table (new series)
from ava.photometry.apertures import ApertureTable
print(ApertureTable(mjd=61000.0))
\end{lstlisting}

\section{Filter and Colour-term Reference}
\label{app:filters}

\begin{table}[h!]
\centering
\small
\begin{tabular}{lllll}
\toprule
LCO filter & GSPC source & APASS band (fallback) & Vega$\to$AB & Typical $K$ \\
\midrule
$u'$ & SDSS $u$ ($-0.04$) & \texttt{apass\_u} & 0.00 & 0.21 \\
$g'$ & SDSS $g$ & \texttt{apass\_g} & 0.00 & 0.14 \\
$r'$ & SDSS $r$ & \texttt{apass\_r} & 0.00 & 0.10 \\
$i'$ & SDSS $i$ & \texttt{apass\_i} & 0.00 & 0.07 \\
$z_s$ & SDSS $z$ ($+0.02$) & --- & 0.00 & 0.07 \\
$B$ & JKC $B$ & \texttt{apass\_B} & $-0.09$ & 0.20 \\
$V$ & JKC $V$ & \texttt{apass\_V} & $+0.02$ & 0.14 \\
$Y$ & PS1 $y$ & --- & 0.00 & 0.07 \\
\bottomrule
\end{tabular}
\caption{LCO filter to reference-catalogue column mapping.
  Extinction coefficients are typical values; actual values per site and dome
  are in \texttt{photometry.extinction\_coefficients} in the config.}
\end{table}

%%============================================================
\section{Legacy Pipeline Reference (v1.3)}
\label{app:legacy}
%%============================================================

\begin{notebox}
The sections below are verbatim extracts from the original legacy pipeline
documentation (\texttt{ava\_pipeline\_manual.tex}).  They describe
\texttt{ava\_update.py} and its child scripts as they existed in production
through 2026, and serve as the specification for the AVA~v2 port.
\end{notebox}


\section{Stage 0 --- Orchestration: \texttt{ava\_update.py}}
\label{sec:orchestrator}
%%============================================================

\texttt{ava\_update.py} is the master driver.  It reads a plain-text target list, loops over
every target, writes one shell script per target that chains the requested pipeline stages,
and launches all scripts simultaneously using \texttt{nohup ... \&}.

\subsection{Target list}

A \textsc{csv} file (default \texttt{target\_list.txt}) with at minimum the columns:
\begin{itemize}[noitemsep]
  \item \texttt{targets} --- the LCO target-name string used in archive queries.
  \item \texttt{directory} --- the per-target folder slug, e.g.\ \texttt{Fairall\_9}.
  \item \texttt{LCO} / \texttt{LT} --- boolean flags enabling the LCO and LT branches.
\end{itemize}

\subsection{Command-line flags}

\begin{center}
\begin{tabular}{cl}
\toprule
Flag & Stage \\
\midrule
\texttt{-o} & Initialise directories (\texttt{ava\_initialise.py}) \\
\texttt{-d} & Download new data from LCO archive \\
\texttt{-e} & Source extraction with SExtractor \\
\texttt{-c} & Photometric calibration \\
\texttt{-p} & Web plots (\texttt{ava\_webplot.py}) \\
\texttt{-m} & Inter-calibration with CALI \\
\texttt{-s} & Save CALI merged output \\
\texttt{-i} & Colour image generation \\
\texttt{-r} & Remove \texttt{lco\_latest*} scratch files \\
\bottomrule
\end{tabular}
\end{center}

The \texttt{-project} argument selects the LCO key-project semester: \texttt{KEY2014A-002},
\texttt{KEY2018B-001}, \texttt{KEY2020B-006}, or \texttt{latest} (auto-detects the most
recent epoch folder present on disk and sets that night as the start date).

\subsection{Script generation}

For each target the driver opens \texttt{1script2rule\textit{dir}.sh}, writes the requested
stage commands as lines, closes the file, makes it executable, and launches it with
\texttt{nohup}.  Log output goes to \texttt{lco\_\textit{dir}\_logall.out}.

\begin{warnbox}
\textbf{Migration target.}  All scripts contain absolute paths hardcoded to
\texttt{/data/jvhs1/agn/} and \texttt{/data3/jvhs1/agn/}.  These must be replaced by
paths read from \texttt{config/ava\_config.yaml} in AVA~v2.
\end{warnbox}

%%============================================================
\section{Stage 1 --- Initialisation: \texttt{ava\_initialise.py}}
\label{sec:init}
%%============================================================

Run once per new target (\texttt{-o}).  Creates the directory tree under
\texttt{master\_dir} and fetches reference catalogues.

\subsection{Directory structure}
\begin{lstlisting}[style=shell]
<master_dir>/<obj>/          # working directory
<master_dir>/<obj>/UVOT/     # Swift/UVOT data placeholder
<master_dir>/raw_data/<obj>/ # raw FITS downloads
\end{lstlisting}

\subsection{APASS catalogue download}

The script uses the \texttt{mechanize} Python browser library to query the APASS
(AAVSO Photometric All-Sky Survey) web form for a cone search centred on the target's
RA/Dec.  The resulting CSV (\texttt{apass\_\textit{obj}.csv}) is the primary photometric
reference catalogue used in calibration.

\begin{warnbox}
\textbf{Migration target.}  This relies on a brittle web scrape that breaks on UI
changes.  AVA~v2 should use \texttt{astroquery.vizier.query\_region()} on APASS\,DR9,
or the Pan-STARRS PS1 catalogue via the LCO archive tools.
\end{warnbox}

%%============================================================
\section{Stage 2 --- Download: \texttt{download\_data.py}}
\label{sec:download}
%%============================================================

Queries the LCO science archive API (\texttt{archive-api.lco.global/frames/}) using an
authentication token for BANZAI-reduced frames (\texttt{rlevel=91}) and downloads them
to the target's raw data directory, organised by night (\texttt{YYYYMMDD/}).

\subsection{Date range logic}

When \texttt{-project latest} is set, the script inspects the most recent epoch folder
present on disk and uses that night as the start date, so only new data is fetched.  For
older semesters the date range is hardcoded (e.g.\ KEY2018B-001: 2018-06-01 to
2020-07-01).

\begin{notebox}
\textbf{AVA v2.}  The download step is replaced entirely by \texttt{lcodb sync} +
\texttt{lcodb daily}, which maintains a local metadata database without bulk-downloading
FITS files.  Images are fetched on demand at reduction time via the archive frame URL.
\end{notebox}

%%============================================================
\section{Stage 3 --- Source Extraction: \texttt{ava\_sextractor.py}}
\label{sec:sextractor}
%%============================================================

Runs SExtractor on every BANZAI-reduced FITS frame (\texttt{*e91.fits.fz}) found under
the target's raw data directory.

\subsection{Input preparation}

BANZAI frames are stored as compressed multi-extension FITS (\texttt{.fz}).  The script
opens each file with \texttt{astropy.io.fits}, extracts the \texttt{SCI} extension, writes
it to a temporary single-extension file (\texttt{temp\_sextractor\_file.fits}), and passes
that to SExtractor.  This step is necessary because SExtractor cannot handle
multi-extension FITS natively.

\subsection{SExtractor configuration}

Fixed configuration files (\texttt{default.sex}, \texttt{default.param},
\texttt{default.conv}, \texttt{default.nnw}, \texttt{default.psfex}) are stored in a
template directory and copied to each target's working folder at initialisation.  The key
output is a source catalogue in FITS binary-table format (\texttt{\_cat.fits}) written to
\texttt{sextractor\_ava/}.

\subsection{Apertures}

SExtractor is run with 12 aperture sizes simultaneously (\texttt{MAG\_APER},
\texttt{MAGERR\_APER}), producing a 12-element array per source.  The science aperture
index used in calibration depends on telescope and epoch (see Section~\ref{sec:aper}).

\subsection{Skip logic}

If a \texttt{\_cat.fits} file already exists and is non-empty for a given frame,
SExtractor is not re-run, making the stage incremental.

%%============================================================
\section{Stage 4 --- Photometric Calibration: \texttt{ava\_calibrate.py}}
\label{sec:calibrate}
%%============================================================

The core science stage.  Reads every BANZAI FITS frame and its SExtractor catalogue,
performs absolute aperture photometry on the AGN and on APASS standard stars, fits a
photometric zeropoint, and appends the result to \texttt{lco\_latest.pkl}.

\subsection{Incremental processing}

If \texttt{lco\_latest.pkl} already exists, only frames whose filename (\texttt{flname})
does not appear in it are processed.  The \texttt{-f} flag forces full reprocessing.

\subsection{Coordinate matching}

The AGN position is read from the target database.  SExtractor source positions are matched
using \texttt{SkyCoord.match\_to\_catalog\_sky()}.  Frames where the nearest source is more
than 3\,arcsec from the AGN position are skipped.

\subsection{Extinction and colour correction}

For each frame, site-specific extinction coefficients are read from
\texttt{lco\_ext\_coefficients.csv} (columns: \texttt{Filter}, \texttt{siteid},
\texttt{dome}, \texttt{K}, \texttt{C\_t}, \texttt{Color1}, \texttt{Color2}).
The corrected instrumental magnitude is:
\begin{equation}
  m_{\rm inst,corr} = m_{\rm inst} + 2.5\log_{10}(t_{\rm exp}) - K\,X - C_t\,c
  \label{eq:ext}
\end{equation}
where $X$ is airmass (minus 1), $t_{\rm exp}$ is exposure time in seconds, and $c$ is the
star's colour from APASS.  Both the AGN and each standard star receive this correction
before zeropoint fitting.


\subsection{Aperture correction}


\subsection{Aperture index resolution}
\label{sec:aper-index}

The SExtractor aperture series changed in 2025: two smaller diameters (3\,px
and 5\,px) were prepended to the original 10-aperture series, shifting every
old index by $+2$.  Using a fixed integer index across epochs therefore
samples \emph{different physical apertures} in the two epochs, corrupting the
light curves silently.

The correct solution, implemented in \texttt{ava.photometry.apertures}, is to
compute the index at reduction time from:
\begin{enumerate}
  \item the \emph{target aperture radius} (a configuration parameter,
        default $= 5$\,\ensuremath{^{\prime\prime}}, i.e.\ 10\,\ensuremath{^{\prime\prime}} diameter), and
  \item the \emph{pixel scale} read from the FITS header \texttt{PIXSCALE} key.
\end{enumerate}

\begin{lstlisting}[style=python]
from ava.photometry.apertures import science_aperture_index

# camera from INSTRUME/filename; mjd from MJD-OBS keyword
idx = science_aperture_index('fa14', mjd=61000.0)   # -> 9
\end{lstlisting}

Tables~\ref{tab:aper-new} and \ref{tab:aper-old} list the physical aperture
diameters for every index and camera family.

\begin{table}[h!]
\centering
\small
\begin{tabular}{rr rrr}
\toprule
 & & \multicolumn{3}{c}{Physical diameter (arcsec)} \\
\cmidrule(lr){3-5}
Index & Diam (px) & Sinistro 1\,m & Spectral 2\,m & MuSCAT3 2\,m \\
      &           & (0.389\,\ensuremath{^{\prime\prime}}/px) & (0.304\,\ensuremath{^{\prime\prime}}/px) & (0.270\,\ensuremath{^{\prime\prime}}/px) \\
\midrule
  0 &  3 &  1.17 &  0.91 &  0.81 \\
  1 &  5 &  1.95 &  1.52 &  1.35 \\
  2 &  8 &  3.11 &  2.43 &  2.16 \\
  3 & 11 &  4.28 &  3.34 &  2.97 \\
  4 & 13 &  5.06 &  3.95 &  3.51 \\
  5 & 15 &  5.83 &  4.56 &  4.05 \\
  6 & 18 &  7.00 &  5.47 &  4.86 \\
  7 & 21 &  8.17 &  6.38 &  5.67 \\
  8 & 24 &  9.34 &  7.30 &  6.48 \\
  9 & 27 & 10.50 &  8.21 &  7.29 \\
 10 & 30 & 11.67 &  9.12 &  8.10 \\
 11 & 33 & 12.84 & 10.03 &  8.91 \\
\bottomrule
\end{tabular}
\caption{New aperture series (2025+, 12 apertures).
  \texttt{PHOT\_APERTURES = 3,5,8,11,13,15,18,21,24,27,30,33} in
  \texttt{default.sex}.  Rows in bold indicate the index selected for a
  5\,arcsec-radius (10\,arcsec-diameter) science aperture on each camera.}
\label{tab:aper-new}
\end{table}

\begin{table}[h!]
\centering
\small
\begin{tabular}{rr rrr}
\toprule
 & & \multicolumn{3}{c}{Physical diameter (arcsec)} \\
\cmidrule(lr){3-5}
Index & Diam (px) & Sinistro 1\,m & Spectral 2\,m & MuSCAT3 2\,m \\
      &           & (0.389\,\ensuremath{^{\prime\prime}}/px) & (0.304\,\ensuremath{^{\prime\prime}}/px) & (0.270\,\ensuremath{^{\prime\prime}}/px) \\
\midrule
  0 &  8 &  3.11 &  2.43 &  2.16 \\
  1 & 11 &  4.28 &  3.34 &  2.97 \\
  2 & 13 &  5.06 &  3.95 &  3.51 \\
  3 & 15 &  5.83 &  4.56 &  4.05 \\
  4 & 18 &  7.00 &  5.47 &  4.86 \\
  5 & 21 &  8.17 &  6.38 &  5.67 \\
  6 & 24 &  9.34 &  7.30 &  6.48 \\
  7 & 27 & 10.50 &  8.21 &  7.29 \\
  8 & 30 & 11.67 &  9.12 &  8.10 \\
  9 & 33 & 12.84 & 10.03 &  8.91 \\
\bottomrule
\end{tabular}
\caption{Old aperture series (pre-2025, 10 apertures).
  \texttt{PHOT\_APERTURES = 8,11,13,15,18,21,24,27,30,33} in
  \texttt{default.sex}.  Note: index~7 here (27\,px, 10.5\,arcsec on 1\,m)
  maps to index~9 in the new series.}
\label{tab:aper-old}
\end{table}

\begin{notebox}
\textbf{Science aperture for re-ingestion.}  For the half-million existing
frames, re-compute the science index via
\texttt{science\_aperture\_index(pixscale, mjd)} rather than storing a fixed
integer.  The function automatically selects the old (pre-MJD\,60909) or new
(MJD\,$\geq$\,60909) series and returns the index closest to the configured
target radius (default 5\,arcsec).  Results: Sinistro 1\,m new series
index~9 (27\,px $= 10.50$"), old series index~7 (27\,px $= 10.50$") ---
the \emph{same physical aperture} in both epochs, confirming the fix is
consistent.  Spectral/MuSCAT3 2\,m: new series index~11
(33\,px $= 10.03$"/8.91"); old series index~9 (same physical size).
\end{notebox}


\label{sec:aper}

Because the science aperture is not necessarily the largest, the script computes aperture
corrections: the median magnitude step between adjacent aperture indices is accumulated
outward to give a correction to infinite aperture for each of the 12 apertures.

The science aperture index is not fixed --- it depends on telescope family and epoch:
\begin{center}
\begin{tabular}{lll}
\toprule
Telescope & Condition & Aperture index \\
\midrule
1\,m & default & 7 \\
1\,m & MJD $> 66909$ & 9 \\
2\,m & MJD $\leq 60909$ & 10 \\
2\,m & MJD $> 60909$ & 12 \\
\bottomrule
\end{tabular}
\end{center}

\begin{notebox}
\textbf{Note.}  The physical aperture sizes in arcsec corresponding to each index depend on
the pixel scale.  For 1\,m Sinistro frames ($0.389$\,arcsec/px): index~7 corresponds to a
diameter of approximately 7\,arcsec; index~9 to approximately 9\,arcsec.  The 2026
expansion was introduced to maintain a consistent metric aperture as the LCO pixel scale
configuration changed. (\textit{Confirm exact MJD thresholds and arcsec values with
J.V.H.S.})
\end{notebox}

\subsection{Zeropoint fitting}

The zeropoint is estimated in two independent ways, both stored in the output:

\textbf{Bootstrap mean} (\texttt{zp0}, \texttt{zp0\_err}).  For each aperture, the median
shift (standard $-$ instrumental corrected) is bootstrap-resampled 50~times to give a
robust mean and uncertainty.

\textbf{\texttt{lmfit} fit} (\texttt{zp0\_fit}, \texttt{zp0\_err}).  A linear fit with
slope fixed to~1.0 and sigma-clipping of outliers yields a second independent estimate.

A minimum of 3~standard stars is required.  If fewer are available, or if the FWHM seeing
exceeds 9\,arcsec, the frame is flagged (\texttt{flag=1}) and stored with NaN photometry.

\subsection{Output columns in \texttt{lco\_latest.pkl}}

One row per FITS frame.  Array-valued columns (aperture vectors of length~12) are stored
as Python lists within DataFrame cells.

\begin{center}
\small
\begin{tabular}{lll}
\toprule
Column & Type & Description \\
\midrule
\texttt{flname}      & str    & Full path to BANZAI FITS file \\
\texttt{Object}      & str    & Target name from FITS header \\
\texttt{Filter}      & str    & LCO filter name \\
\texttt{MJD}         & float  & Mid-exposure MJD \\
\texttt{mag\_aper}   & list[12] & Aperture magnitudes (instrumental + offset) \\
\texttt{err\_aper}   & float  & Photometric error at science aperture \\
\texttt{zp\_mag}     & list[12] & Bootstrap zeropoints per aperture \\
\texttt{zp\_err}     & list[12] & Zeropoint uncertainties \\
\texttt{zp0\_fit}    & float  & \texttt{lmfit} zeropoint \\
\texttt{zp0\_err}    & float  & \texttt{lmfit} zeropoint uncertainty \\
\texttt{corrects}    & list[12] & Aperture corrections \\
\texttt{extra\_fair} & float  & Extinction + colour correction applied to AGN \\
\texttt{flag}        & int    & 0 = good; 1 = too few standards or bad seeing \\
\texttt{n\_std}      & int    & Number of standard stars used \\
\texttt{dist}        & float  & AGN--catalogue separation (arcsec) \\
\texttt{seeing}      & float  & Measured FWHM (arcsec) \\
\texttt{airmass}     & float  & Airmass $-1$ \\
\texttt{observatory} & str    & LCO site ID \\
\texttt{telid}       & str    & Telescope ID from filename \\
\texttt{skymag}      & float  & Sky background magnitude \\
\texttt{m\_fit}      & float  & Slope of linear photometric solution \\
\texttt{b\_fit}      & float  & Intercept of linear photometric solution \\
\texttt{X}, \texttt{Y} & float & AGN pixel position in image \\
\bottomrule
\end{tabular}
\end{center}

%%============================================================
\section{Stage 5 --- Plotting \& Web Output: \texttt{ava\_webplot.py}}
\label{sec:webplot}
%%============================================================

Reads \texttt{lco\_latest.pkl} and produces three sets of plots per target plus writes
calibrated light-curve data as text files for the website.

\subsection{Flux conversion}

Calibrated magnitudes are converted to flux density in mJy:
\begin{equation}
  f_\nu \,[\text{mJy}] = 10^{(m + \mathrm{ZP} - 8.9)\,/\,-2.5} \times 10^3
  \label{eq:flux}
\end{equation}
with uncertainties $\sigma_f = f_\nu \,/\, 1.0857 \times \sigma_m$.

\subsection{Quality cuts applied at plot time}

\begin{itemize}[noitemsep]
  \item \texttt{flag == 0} (calibration succeeded)
  \item \texttt{dist < 3.5}\,arcsec (AGN position match)
  \item \texttt{n\_std >= 3} (at least 3~standard stars)
\end{itemize}

\subsection{Plots produced}

\begin{enumerate}
  \item \textbf{Full light curve} (\texttt{lco\_latest\_\textit{obj}.pdf/png}).  One panel
        per filter over the full time baseline.  FLOYDS epoch times are marked as vertical
        ticks on the top panel.
  \item \textbf{Last 60~days} (\texttt{lco\_latest\_\textit{obj}\_last6.pdf/png}).  Same
        layout restricted to the 60~days preceding the most recent observation.  Y-axis
        limits are set from the 4th--96th flux percentile with 50\,\% padding.
  \item \textbf{Swift} (\texttt{lco\_latest\_\textit{obj}\_swift.pdf/png}).  Reads
        \texttt{ava\_swift\_updated.csv} and plots UV and X-ray light curves.
\end{enumerate}

\subsection{Diagnostic plots (\texttt{-d} flag)}

\begin{itemize}[noitemsep]
  \item Seeing, airmass, and zeropoint (transparency) vs time per site.
  \item Standard-star magnitude vs photometric RMS, with the AGN overlaid.  This is the
        primary internal photometric quality diagnostic.
\end{itemize}

\subsection{Remote delivery}

Plots are \texttt{scp}\textquotesingle d to the web server
(\texttt{ringworld.st-andrews.ac.uk:/polaris2/jvhs1/halpha\_survey/ava\_latest/}) by the
generated shell script.

%%============================================================
\section{Stage 6 --- Inter-calibration: \texttt{ava\_merge.py}}
\label{sec:merge}
%%============================================================

Photometry from different telescopes, cameras, and sites cannot be concatenated directly:
each has its own zeropoint offset.  \texttt{ava\_merge.py} uses
\texttt{pycali\_jvhs} --- a custom Python wrapper around the CALI MCMC
inter-calibration algorithm --- to solve for per-telescope flux scale factors and merge
the light curves into a single consistent time series.

\subsection{Input}

The script reads \texttt{<obj>\_edited.csv} from the parent directory.  This is the
calibrated light-curve file (essentially \texttt{lco\_latest.pkl} converted to CSV and
optionally hand-edited to remove bad epochs).  It must exist before merging can proceed.

\subsection{CALI algorithm}

For each filter, unique telescope identifiers (\texttt{telid}) are found.
\texttt{pycali\_jvhs.pycali()} is run with an MCMC chain of $10{,}000$ steps (extended
by $10{,}000$ on each retry if convergence is not achieved).  CALI models the light curve
as a common signal observed by each telescope with an additive offset and a flux scaling
factor, marginalising over these nuisance parameters.

\subsection{Output}

Merged light-curve files and plots are written to the \texttt{ava\_merging/} subdirectory,
then copied to \texttt{ava\_plots/} and \texttt{scp}\textquotesingle d to the web server.

\begin{warnbox}
\textbf{Migration target.}  \texttt{pycali\_jvhs} is a private package not on PyPI and
must be installed separately.  AVA~v2 should document the install procedure or replace
it with PyROA or CREAM.
\end{warnbox}

%%============================================================
\section{Data Products \& File Conventions}
\label{sec:products}
%%============================================================

All intermediate and final products are stored per-target under a fixed directory tree
rooted at \texttt{master\_dir} (currently \texttt{/data3/jvhs1/agn/}).

\begin{longtable}{>{\ttfamily}p{5.5cm}p{8.5cm}}
\toprule
File / folder & Contents \\
\midrule
apass\_\textit{obj}.csv & APASS reference star catalogue \\
apass\_pytics.csv & APASS catalogue reformatted for pytics column names \\
lco\_latest.pkl & Running pandas DataFrame of all calibrated photometry \\
lco\_latest\_stan.pkl & Running DataFrame of all calibrated standard stars \\
lco\_latest\_\textit{obj}.pdf/png & Full light-curve plot \\
lco\_latest\_\textit{obj}\_last6.pdf/png & Last-60-days light-curve plot \\
lco\_latest\_\textit{obj}\_swift.pdf/png & Swift UV + X-ray plot \\
lco\_latest\_\textit{obj}\_\textit{filt}\_sky.pdf/png & Diagnostic plots \\
lco\_latest\_\textit{obj}\_stan.pdf/png & Standard-star RMS plot \\
sextractor\_ava/ & SExtractor catalogues (\texttt{*\_cat.fits}) \\
ava\_merging/ & CALI workspace and merged output files \\
UVOT/ & Swift/UVOT data placeholder \\
default.sex, etc. & SExtractor configuration files \\
filter\_dict.txt & Filter name / wavelength / colour mapping \\
lco\_ext\_coefficients.csv & Extinction and colour-term coefficients \\
\bottomrule
\caption{Per-target working directory contents.}
\end{longtable}

\textbf{Raw FITS data} are stored at
\texttt{<master\_dir>/raw\_data/<obj>/YYYYMMDD/*e91.fits.fz}, one sub-folder per
observing night.

\begin{warnbox}
\textbf{Pickle portability.}  \texttt{lco\_latest.pkl} stores 12-element aperture
vectors as Python lists inside a pandas DataFrame.  Reading them requires
\texttt{np.stack(lco.mag\_aper[mask], axis=0)}.  This format is fragile across
pandas/numpy versions.  AVA~v2 should migrate to HDF5 or Parquet, or store photometry
directly in the \texttt{lcodb} SQLite database.
\end{warnbox}

%%============================================================

\end{document}
